\subsection{Restrições de Segurança e Zonas Proibidas (Keep-Out Zones)}

Além das restrições apresentadas anteriormente, é necessário considerar as restrições geométricas associadas ao entorno do ponto de captura.  
Essas zonas de exclusão, também chamadas de \emph{Keep-Out Zones} (KOZ), correspondem a regiões ao redor do alvo em que o manipulador robótico não deve ingressar, sob risco de colisão ou de interferência mecânica.

A restrição pode ser formulada como uma condição de exclusão espacial:
\begin{equation}
\label{eq:koz}
\vec{p}_E(t) \notin \mathcal{Z}_{\text{KOZ}} \quad \forall t \in [t_0, t_{\text{captura}}],
\end{equation}

onde:
\begin{itemize}
    \item $\vec{p}_E(t)$: posição do efetuador ao longo do tempo;
    \item $\mathcal{Z}_{\text{KOZ}}$: região tridimensional proibida (zona de exclusão);
    \item $t_0$: início da aproximação fina;
    \item $t_{\text{captura}}$: instante de captura do alvo.
\end{itemize}

A forma e o tamanho da zona $\mathcal{Z}_{\text{KOZ}}$ dependem da geometria do sistema de acoplamento, das tolerâncias de aproximação e do tipo de missão.  
Tais zonas são definidas no planejamento da missão e devem ser respeitadas pelo sistema de controle do manipulador.

Além da restrição espacial (KOZ), é necessário garantir que o alvo permaneça dentro do campo de visão do sensor de bordo durante toda a aproximação.  
Essa condição pode ser descrita pela restrição angular entre o eixo óptico do sensor ($\vec{b}$) e o vetor em direção ao alvo ($\vec{r}$):

\begin{equation}
\label{eq:fov}
\frac{\vec{b}^T \vec{r}}{\|\vec{r}\|} \geq \cos \varphi_{\min},
\end{equation}

onde:
\begin{itemize}
    \item $\vec{b}$: vetor unitário que define o eixo óptico do sensor;
    \item $\vec{r}$: vetor que aponta para o alvo;
    \item $\varphi_{\min}$: ângulo mínimo admissível para manter o alvo no campo de visão.
\end{itemize}

Para garantir conformidade com essas restrições, métodos de navegação segura e verificação de colisão são incorporados ao sistema de controle, inclusive durante simulações em ambiente de microgravidade.
